\documentclass[11pt,a4paper]{article}

\usepackage[utf8]{inputenc}
\usepackage[T1]{fontenc}
\usepackage[spanish]{babel}
\usepackage{geometry}
\usepackage{xcolor}
\usepackage{booktabs}
\usepackage{longtable}
\usepackage{tabularx}
\usepackage{array}
\usepackage{enumitem}
\usepackage{hyperref}
\usepackage{fancyhdr}
\usepackage{listings}
\usepackage{graphicx}

\geometry{margin=2.1cm}
\definecolor{azulmanzanares}{HTML}{126D98}
\definecolor{verdemanzanares}{HTML}{168052}
\definecolor{grisclaro}{HTML}{F2F5F7}
\definecolor{grisoscuro}{HTML}{38444D}
\hypersetup{
  colorlinks=true,
  linkcolor=azulmanzanares,
  urlcolor=azulmanzanares,
  pdftitle={MANZANARES Growth Agent - Propuesta técnica de piloto},
  pdfauthor={Documento de trabajo para Grupo Manzanares S.A.S.}
}
\pagestyle{fancy}
\setlength{\headheight}{14pt}
\fancyhf{}
\lhead{MANZANARES Growth Agent}
\rhead{Propuesta técnica de piloto}
\cfoot{\thepage}
\setlist{itemsep=2pt,topsep=4pt}
\renewcommand{\arraystretch}{1.15}

\lstset{
  basicstyle=\ttfamily\small,
  backgroundcolor=\color{grisclaro},
  frame=single,
  breaklines=true,
  columns=fullflexible,
  showstringspaces=false
}

\newcommand{\decisionbox}[1]{
  \begin{center}
  \fcolorbox{azulmanzanares}{grisclaro}{
    \parbox{0.91\textwidth}{\textbf{Decisión ejecutiva:} #1}
  }
  \end{center}
}

\begin{document}
\hypersetup{pageanchor=false}

\begin{titlepage}
\centering
\vspace*{1.6cm}
{\color{azulmanzanares}\rule{\textwidth}{2.5pt}}\\[1.2cm]
{\Huge\bfseries MANZANARES Growth Agent\par}
\vspace{0.35cm}
{\Large Propuesta técnica de piloto empresarial\par}
\vspace{0.25cm}
{\large Contact center de la UEN Carnes Manzanares\par}
\vspace{1.4cm}

\fcolorbox{verdemanzanares}{grisclaro}{
  \parbox{0.82\textwidth}{
    \centering
    \textbf{Objetivo}\\[4pt]
    Convertir la gestión comercial en un proceso medible, trazable y escalable,
    capaz de aumentar ventas mediante datos, priorización y disciplina de
    seguimiento, sin crecer proporcionalmente el equipo operativo.
  }
}

\vfill
\begin{tabular}{rl}
\textbf{Organización:} & Grupo Manzanares S.A.S.\\
\textbf{Sector:} & Agroalimentos\\
\textbf{Duración propuesta:} & 16 semanas\\
\textbf{Presupuesto de referencia:} & 21.000.000 COP\\
\textbf{Versión de solución:} & 2.0 - núcleo auditable de piloto\\
\textbf{Fecha:} & 25 de junio de 2026\\
\end{tabular}
\vfill
{\color{azulmanzanares}\rule{\textwidth}{2.5pt}}
\end{titlepage}

\hypersetup{pageanchor=true}
\pagenumbering{roman}
\tableofcontents
\clearpage
\pagenumbering{arabic}

\section{Resumen ejecutivo}

La oportunidad no consiste únicamente en incorporar inteligencia artificial. El
valor principal está en instalar un sistema operativo comercial: datos mínimos
confiables, una cola diaria priorizada, seguimiento obligatorio, experimentos
medibles y una lectura ejecutiva común.

Se propone un piloto de 16 semanas que combina un CRM mínimo, agentes
especializados, scoring explicable, auditoría y medición experimental. La versión
2.0 implementada ya incluye migraciones de datos, controles de integridad, API
interna, autenticación opcional, health checks, backups, importación CSV validada,
operación sin dependencias externas y pruebas automatizadas.

El piloto no se presenta como un CRM corporativo terminado. Se presenta como un
mecanismo de validación de impacto y adopción que deja una base escalable hacia
PostgreSQL, identidad corporativa e integraciones empresariales.

\decisionbox{Aprobar un piloto controlado orientado a evidencia. La decisión de
escalamiento debe depender de adopción, calidad de datos e ingreso incremental,
no de la sofisticación aparente de la tecnología.}

\section{Entendimiento del reto}

\subsection{Situación actual}

La ficha describe una operación principalmente manual, subjetiva y reactiva, con
limitaciones en prospección, seguimiento, recuperación de clientes inactivos y
recompra. No existe trazabilidad completa del ciclo
datos-cliente-contacto-oportunidad-pedido-recompra.

Las áreas involucradas incluyen dirección comercial, mercadeo, contact center,
líder de la UEN, puntos de venta y tecnología. La solución debe convivir con
WhatsApp, llamadas y la operación vigente, manteniendo simplicidad de uso.

\subsection{Meta comercial por aclarar}

La ficha utiliza dos formulaciones incompatibles:

\begin{enumerate}
  \item duplicar las ventas actuales; y
  \item aumentar 50\%, pasando de 40 a 80 millones COP.
\end{enumerate}

Pasar de 40 a 80 millones representa un aumento de 100\%. La meta contractual,
la venta atribuible al contact center y la línea base deben quedar aprobadas en
la fase de descubrimiento.

\section{Propuesta de valor}

MANZANARES Growth Agent convierte datos comerciales en decisiones diarias:

\begin{enumerate}
  \item consolida contactos, oportunidades, interacciones y tareas;
  \item calcula línea base, pipeline, brecha y calidad de datos;
  \item segmenta clientes y prospectos por estado, valor y recurrencia;
  \item asigna un score explicable y muestra las razones;
  \item crea tareas persistentes sin duplicar ni borrar historial;
  \item registra resultados y alimenta el siguiente ciclo;
  \item diseña pilotos con métricas y comparación válida;
  \item produce un briefing ejecutivo determinístico o mejorado localmente.
\end{enumerate}

\subsection{Experiencia diaria esperada}

\begin{longtable}{p{3cm}p{5.2cm}p{6.4cm}}
\toprule
\textbf{Momento} & \textbf{Usuario} & \textbf{Experiencia}\\
\midrule
\endhead
Inicio de jornada & Líder comercial & Verifica salud, tareas vencidas, brecha y
distribución por asesor\\
Priorización & Asesor & Recibe cola ordenada con razón de prioridad y siguiente
acción\\
Gestión & Asesor & Registra canal, resultado, nota y próxima acción\\
Cierre diario & Líder & Revisa cobertura, completitud y pendientes\\
Revisión semanal & Comité del piloto & Evalúa embudo, conversión, ingreso
incremental, adopción y riesgos\\
\bottomrule
\end{longtable}

\section{Arquitectura de la solución}

\subsection{Capas}

\begin{center}
\fbox{\parbox{0.9\textwidth}{
\centering
\textbf{CLI y API interna}\\
$\downarrow$\\
\textbf{Coordinador multi-intención}\\
$\downarrow$\\
\textbf{Diagnóstico - Segmentación - Priorización - Seguimiento -
Reactivación - Growth}\\
$\downarrow$\\
\textbf{CRM versionado - Auditoría - Health checks - Backups}\\
$\downarrow$\\
\textbf{Síntesis determinística - Ollama local opcional}
}}
\end{center}

\subsection{Agentes}

\begin{longtable}{p{2.8cm}p{5cm}p{6.8cm}}
\toprule
\textbf{Agente} & \textbf{Responsabilidad} & \textbf{Salida}\\
\midrule
\endhead
Coordinador & Detectar una o varias intenciones & Flujo combinado de agentes\\
Diagnóstico & Cuantificar línea base, brecha y capacidad & KPIs, completitud,
actividad y alertas\\
Segmentador & Clasificar la base & Grupos accionables y registros incompletos\\
Priorizador & Calcular score reproducible & Ranking, componentes y razones\\
Seguimiento & Sincronizar la cola diaria & Tareas nuevas, actualizadas,
reemplazadas y conservadas\\
Reactivación & Identificar recuperación potencial & Cohorte, valor histórico y
diseño controlado\\
Growth & Diseñar experimentos & Hipótesis, métrica primaria y regla de parada\\
Redactor & Consolidar evidencia & Briefing ejecutivo sin inventar cifras\\
\bottomrule
\end{longtable}

\subsection{Decisiones tecnológicas}

\begin{itemize}
  \item Python 3.11 o superior, sin dependencias de runtime.
  \item SQLite para piloto de baja concurrencia, con WAL, claves foráneas,
  índices, timeout y backups consistentes.
  \item API HTTP interna para integración progresiva.
  \item Contenedor no privilegiado, filesystem de solo lectura y health check.
  \item LLM local desactivado por defecto; el núcleo comercial es determinístico.
\end{itemize}

\section{Controles de nivel empresarial incorporados}

\begin{longtable}{p{4cm}p{5.2cm}p{5.4cm}}
\toprule
\textbf{Brecha del prototipo} & \textbf{Mejora implementada} & \textbf{Resultado}\\
\midrule
\endhead
Esquema sin evolución & Migraciones versionadas y detección de base heredada &
Actualización sin reinstalación manual\\
Claves foráneas desactivadas & Integridad referencial por conexión & Menor riesgo
de registros huérfanos\\
Bloqueo y concurrencia básica & WAL, timeout e índices & Mejor respuesta del
piloto local\\
Tareas borradas en cada ejecución & Sincronización idempotente y estado
\texttt{superseded} & Historial preservado y cola sin duplicados\\
Prospectos con recencia ficticia de 999 días & Score corregido por estado &
Priorización más coherente\\
Una sola intención por consulta & Orquestación multi-intención & Respuestas que
combinan brecha, prioridades y campañas\\
Auditoría con texto sensible & Hash de pregunta y metadatos operativos &
Trazabilidad con menor exposición\\
LLM como dependencia implícita & Fallback determinístico, timeout y límites &
Continuidad operativa\\
Una sola prueba no descubierta & Suite automatizada y descubrible & Evidencia de
regresión y calidad\\
Sin interfaz de integración & API interna y salida JSON & Preparación para
canales corporativos\\
\bottomrule
\end{longtable}

\section{Modelo de datos y gobierno}

El CRM mínimo utiliza:

\begin{itemize}
  \item \texttt{contacts}: identidad externa, segmento, estado, potencial,
  recurrencia, interés, responsable, consentimiento y fuente;
  \item \texttt{opportunities}: etapa, valor, probabilidad y siguiente acción;
  \item \texttt{interactions}: canal, resultado, nota y fecha;
  \item \texttt{tasks}: prioridad, vencimiento, estado, razón y ejecución origen;
  \item \texttt{system\_runs}: estado, duración, agentes y uso de LLM;
  \item \texttt{audit\_log}: eventos técnicos y métricas por agente.
\end{itemize}

Antes de utilizar datos reales deben aprobarse propietario, custodio, finalidad,
clasificación, autorización, retención, eliminación y separación de ambientes.
Los datos de demostración deben permanecer sintéticos o enmascarados.

\section{Scoring comercial explicable}

\begin{center}
\begin{tabular}{lr}
\toprule
\textbf{Dimensión} & \textbf{Peso máximo}\\
\midrule
Potencial económico & 30\\
Interés observado & 25\\
Tipo de relación & 15\\
Urgencia o siguiente acción & 15\\
Frecuencia histórica & 10\\
Calidad mínima del registro & 5\\
\midrule
\textbf{Total} & \textbf{100}\\
\bottomrule
\end{tabular}
\end{center}

El score es una ayuda para ordenar trabajo, no una decisión automática de
exclusión. Cada resultado muestra componentes y razones. Los pesos deben
recalibrarse con pedidos y conversión real.

\section{Seguridad, privacidad y continuidad}

\subsection{Controles del piloto}

\begin{itemize}
  \item API limitada a localhost por defecto.
  \item Token obligatorio al escuchar en interfaces externas.
  \item Límite de solicitudes y cabeceras HTTP de endurecimiento.
  \item Secretos fuera del repositorio.
  \item Auditoría por ejecución y degradación controlada del LLM.
  \item Usuario no privilegiado en contenedor.
  \item Backups consistentes mediante API de SQLite.
\end{itemize}

\subsection{Requisitos antes de producción}

\begin{itemize}
  \item identidad corporativa, roles y proceso de altas y bajas;
  \item TLS, cifrado administrado y rotación de secretos;
  \item logs, métricas, alertas y backups centralizados;
  \item pruebas de restauración, vulnerabilidades y penetración;
  \item política aprobada de tratamiento y conservación de datos;
  \item migración a PostgreSQL cuando la concurrencia o criticidad lo requiera.
\end{itemize}

\section{Indicadores del piloto}

\subsection{Embudo y valor}

\begin{itemize}
  \item contactos asignados, intentos y contactos efectivos;
  \item oportunidades, propuestas, pedidos y conversión por etapa;
  \item pipeline bruto y ponderado;
  \item ingreso atribuible e ingreso incremental;
  \item brecha frente a línea base y meta aprobadas.
\end{itemize}

\subsection{Retención y productividad}

\begin{itemize}
  \item clientes con 30, 60 y 90 días sin compra;
  \item tasa de reactivación y recompra;
  \item tareas vencidas y tiempo hasta la siguiente acción;
  \item acciones completadas por asesor;
  \item completitud de datos y adopción semanal.
\end{itemize}

\section{Pilotos recomendados}

\subsection{Piloto A: reactivación}

Seleccionar clientes con más de 60 días sin compra y asignarlos a tratamiento y
control. Medir contacto efectivo, pedido, ingreso incremental y recompra.

\subsection{Piloto B: prospección institucional}

Priorizar restaurantes, hoteles, catering y comedores por potencial, necesidad y
probabilidad. Medir contactos, reuniones, muestras, propuestas y pedidos.

\subsection{Piloto C: recompra programada}

Generar tareas antes del ciclo esperado de compra y comparar mensajes por
segmento. Medir recompra, anticipación y valor por cliente.

\section{Plan de 16 semanas}

\begin{longtable}{p{2cm}p{4.5cm}p{8cm}}
\toprule
\textbf{Semanas} & \textbf{Fase} & \textbf{Entregables y puerta de control}\\
\midrule
\endhead
1 a 2 & Descubrimiento & Meta aprobada, línea base, proceso, fuentes, riesgos,
diccionario y criterios de éxito\\
3 a 4 & Datos y CRM mínimo & Importación, calidad, responsables, permisos,
migraciones y backup probado\\
5 a 6 & Priorización & Segmentos, score, cola diaria, API y capacitación inicial\\
7 a 9 & Piloto A & Reactivación con tratamiento, control y medición incremental\\
10 a 12 & Pilotos B y C & Prospección institucional y recompra programada\\
13 a 14 & Calibración & Pesos, mensajes, automatizaciones y experiencia operativa\\
15 & Transferencia & Manuales, seguridad, gobierno, restauración y soporte\\
16 & Cierre & Resultados, aceptación, backlog y decisión de escalamiento\\
\bottomrule
\end{longtable}

\section{Presupuesto de referencia}

La siguiente distribución es una propuesta de planeación y debe ajustarse con el
alcance definitivo:

\begin{center}
\begin{tabularx}{0.92\textwidth}{Xr}
\toprule
\textbf{Frente} & \textbf{COP}\\
\midrule
Descubrimiento, proceso y gobierno & 3.000.000\\
Datos, CRM mínimo e importación & 5.000.000\\
Agentes, API y controles operativos & 4.500.000\\
Pilotos y medición experimental & 4.000.000\\
Capacitación y transferencia & 2.000.000\\
Contingencia, cierre y ruta de escalamiento & 2.500.000\\
\midrule
\textbf{Total} & \textbf{21.000.000}\\
\bottomrule
\end{tabularx}
\end{center}

No se recomienda comprometer dentro de este presupuesto un CRM empresarial
completo, integración omnicanal profunda, alta disponibilidad o modelos
predictivos entrenados antes de validar calidad y volumen de datos.

\section{Criterios de aceptación}

\begin{itemize}
  \item al menos 95\% de registros priorizados con teléfono, segmento y asesor;
  \item cola diaria sin duplicados y con historial de reemplazos;
  \item interacciones trazables y tareas cerrables;
  \item health check correcto y restauración demostrada;
  \item score con razones visibles;
  \item definición aprobada de venta atribuible, línea base y meta;
  \item evidencia de mejora frente a línea base o grupo de comparación;
  \item plan aprobado para seguridad e integración antes de producción.
\end{itemize}

\section{Riesgos y controles}

\begin{longtable}{p{3.8cm}p{5cm}p{5.8cm}}
\toprule
\textbf{Riesgo} & \textbf{Impacto} & \textbf{Control}\\
\midrule
\endhead
Datos incompletos & Priorización y medición débiles & Perfilado, campos mínimos,
rechazo de errores y responsable por fuente\\
Meta ambigua & Conflicto de evaluación & Aprobación formal en semanas 1 a 2\\
Baja adopción & Herramienta sin impacto & Diseño con asesores, capacitación y
tablero de uso\\
Inventario desactualizado & Ofertas inviables & Validación de vigencia antes de
campañas\\
Datos personales & Riesgo legal y reputacional & Minimización, roles, retención,
auditoría y controles corporativos\\
Dependencia de IA & Interrupción o respuestas no confiables & Operación
determinística y LLM opcional\\
Expectativa de producto final & Sobrecosto y retraso & Alcance de piloto,
exclusiones y puertas de decisión\\
SQLite fuera de contexto & Concurrencia insuficiente & Umbrales de migración a
PostgreSQL\\
\bottomrule
\end{longtable}

\section{Evidencia de implementación}

El repositorio entrega actualmente:

\begin{itemize}
  \item CLI y API interna;
  \item migración automática desde la base inicial;
  \item CRM SQLite con integridad, WAL e índices;
  \item seis agentes funcionales y coordinador multi-intención;
  \item scoring explicable corregido;
  \item tareas idempotentes con historial;
  \item importación CSV con validación previa;
  \item auditoría por ejecución;
  \item health checks y backups;
  \item contenedor endurecido;
  \item suite automatizada de 15 pruebas;
  \item documentación de seguridad, operación, datos y aceptación.
\end{itemize}

\subsection{Comandos de verificación}

\begin{lstlisting}
python3 app.py inicializar
python3 app.py salud
python3 app.py tablero
python3 app.py consultar "Dame la brecha y las prioridades" --json
python3 app.py importar-contactos data/contactos_ejemplo.csv --dry-run
python3 -m unittest discover -v
\end{lstlisting}

\section{Información requerida a Grupo Manzanares}

\begin{enumerate}
  \item histórico de clientes, pedidos, productos, fechas, cantidades y valores;
  \item estados actuales de prospectos y oportunidades;
  \item registros disponibles de llamadas y WhatsApp;
  \item asesores, zonas, segmentos y reglas de asignación;
  \item catálogo e inventario con frecuencia de actualización;
  \item definición de venta atribuible al contact center;
  \item línea base validada y estacionalidad;
  \item política y autorizaciones de tratamiento de datos;
  \item arquitectura actual y posibilidades de integración;
  \item confirmación formal de la meta comercial.
\end{enumerate}

\section{Recomendación final}

La solución ya demuestra el flujo esencial, pero el éxito empresarial dependerá
de tres factores: calidad de datos, disciplina de registro y diseño experimental.
La recomendación es iniciar el piloto con una cohorte acotada, una meta aprobada
y revisiones semanales. El escalamiento debe aprobarse únicamente cuando exista
evidencia conjunta de impacto comercial, adopción operativa y control de riesgo.

\subsection{Decisión solicitada al comité}

\begin{enumerate}
  \item autorizar la fase de descubrimiento y acceso controlado a fuentes;
  \item designar un dueño de negocio, un custodio de datos y representantes de
  asesores;
  \item aprobar la puerta de decisión al cierre de la semana 2, cuando estén
  validadas meta, línea base, alcance, riesgos y criterios de éxito.
\end{enumerate}

\subsection{Primeros diez días hábiles}

\begin{center}
\begin{tabularx}{0.94\textwidth}{p{2.3cm}X}
\toprule
\textbf{Momento} & \textbf{Actividad}\\
\midrule
Días 1 y 2 & Alineación ejecutiva, responsables, alcance, restricciones y
definición preliminar de éxito\\
Días 3 a 5 & Perfilado de fuentes, diccionario, calidad, permisos y mapa del
proceso actual\\
Días 6 a 8 & Línea base, definición de venta atribuible, segmentos y cohorte
inicial\\
Días 9 y 10 & Presentación de hallazgos, riesgos, diseño del piloto y decisión de
continuidad\\
\bottomrule
\end{tabularx}
\end{center}

\decisionbox{El entregable de la semana 2 debe permitir decidir con evidencia si
el piloto puede alcanzar un impacto medible dentro de 16 semanas y 21 millones
COP, o si requiere ajustar meta, alcance o fuentes antes de continuar.}

\end{document}
